% =====================================================================
%  Anexos (antes de las Referencias, según la plantilla ITBA)
% =====================================================================
\chapter*{Anexos}
\addcontentsline{toc}{chapter}{Anexos}
\addtocontents{toc}{\protect\anexotocspacing}
\setcounter{section}{0}
\renewcommand{\thesection}{Anexo \Alph{section}}
\setcounter{table}{0}
\renewcommand{\thetable}{\Alph{section}.\arabic{table}}
\setcounter{figure}{0}
\renewcommand{\thefigure}{\Alph{section}.\arabic{figure}}

\section{Tamaño muestral por grupo en la prueba SLR}\label{anx:muestra}

\begin{table}[htbp]
\centering
\caption[Número de animales por grupo en la prueba SLR]{Número de animales por grupo experimental (edad de exposición $\times$ tratamiento) en cada configuración de la prueba SLR. Algunos animales realizaron ambas configuraciones, por lo que las columnas no son aditivas.}
\label{tab:n-slr}
\small
\begin{tabular}{@{}lcc@{}}
\toprule
\textbf{Grupo} & \textbf{d-SLR} & \textbf{s-SLR} \\
\midrule
Juvenil -- Control  & 9  & 10 \\
Juvenil -- Sacarosa & 9  & 9  \\
Juvenil -- Stevia   & 10 & 10 \\
Adulto -- Control   & 10 & 10 \\
Adulto -- Sacarosa  & 10 & 10 \\
Adulto -- Stevia    & 10 & 10 \\
\bottomrule
\end{tabular}

\vspace{0.4em}
\begin{minipage}{0.9\linewidth}
\footnotesize\textit{Los grupos control provienen exclusivamente de las cohortes 2025; ninguna cohorte de 2026 incluyó grupo control propio.}
\end{minipage}
\end{table}
